\documentclass{article}
\usepackage[margin=2.5cm]{geometry} % Riduce i margini in alto, in basso e ai lati
\usepackage{graphicx} % Per l'inserimento delle immagini
\usepackage{booktabs} % Per la formattazione professionale delle tabelle (\toprule, \midrule, \bottomrule)
\usepackage{xurl}     % Gestione avanzata dei punti di a capo
\usepackage{parskip}

\title{Progetto Deep Learning}
\author{Lorenzo Didomenico}
\date{Ottobre 2026}

\begin{document}

\maketitle
 \clearpage
% --- AGGIUNGI DA QUI ---
\tableofcontents
\clearpage % Sposta l'Introduzione nella pagina successiva
% --- A QUI ---

\section{Introduzione}

I Video Large Language Model hanno raggiunto buoni risultati su video brevi e
con un contenuto visivo esplicito, ma restano deboli quando la risposta richiede
di seguire un'attività lunga e ripresa in prima persona. HD-EPIC~\cite{hdepic} è
un dataset egocentrico di attività di cucina annotato in modo molto dettagliato,
a cui è associato un benchmark VQA a scelta multipla costruito proprio per
testare questi modelli: nel paper originale i modelli testati restano molto
lontani dall'accuratezza degli annotatori umani.

Questo progetto si concentra su una sola categoria del benchmark, \emph{Nutrition},
composta da 200 domande a cinque opzioni divise in tre prototipi. È la categoria
in cui la risposta dipende dall'individuare gli ingredienti lungo la preparazione ed in particolare dare risposte riguardo i loro valori nutrizionali.

Gli obiettivi sono due. Il primo è \textbf{riprodurre le baseline} riportate in
HD-EPIC (Tab.~2, colonna \emph{Nutrition}) per due modelli open da 7B,
VideoLLaMA~2~\cite{videollama2} e LLaVA-Video~\cite{llavavideo}, per cui il
paper riporta rispettivamente 32.7\% e 38.7\% contro un 20\% casuale. Il secondo
è \textbf{valutare un metodo di visual prompting}, i Video
Panels~\cite{videopanels}, che a parità di finestra di contesto sostituisce ogni
frame con una griglia $2\times2$ di frame consecutivi: si quadruplica
così la copertura temporale al prezzo di un quarto della risoluzione spaziale.

Entrambi gli esperimenti passano da un unico modulo di valutazione, condiviso
dai due modelli che fissa il prompt, il parsing delle risposte e le metriche:
l'unica parte specifica del modello è la funzione di inferenza. Questo è ciò che
rende confrontabili i risultati, perché con parser diversi la differenza di
accuracy fra due modelli misurerebbe anche la differenza fra i parser.

I risultati ottenuti sono i seguenti. Le baseline sono riprodotte su tutte le
200 domande: VideoLLaMA~2 ottiene $36.7\%$, $4.0$ punti sopra il valore
pubblicato, mentre LLaVA-Video ottiene $38.7\%$, esattamente il valore
pubblicato. I Video Panels, invece, non producono alcun guadagno
sulla categoria: $0.0$ punti per VideoLLaMA~2 e $-1.0$ per LLaVA-Video. Il dato
più interessante non è però il risultato aggregato, ma il fatto che i due
modelli si spostino in direzioni opposte sui due prototipi: VideoLLaMA~2 guadagna con i Video Panels su \texttt{nutrition\_change} e perde su \texttt{video\_nutrition\_estimation}, mentre LLaVA-Video si comporta in maniera opposta.

\section{Dataset e benchmark}

HD-EPIC~\cite{hdepic} raccoglie 41 ore di video egocentrici
registrati in 9 cucine domestiche, per un totale di 156 video di durata
media pari a 15.9 minuti. La caratteristica del dataset è il numero delle
annotazioni, in media 263 al minuto, che coprono azioni fine-grained, passaggi
delle ricette, eventi audio, movimenti degli oggetti, digital twin 3D delle
cucine e direzione dello sguardo.

Su queste annotazioni è costruito un benchmark VQA a scelta multipla di circa
26.000 domande, suddivise in sette categorie: \emph{Recipes}, \emph{Ingredients},
\emph{Nutrition}, \emph{Fine-grained Actions}, \emph{3D Perception},
\emph{Object Motion} e \emph{Gaze}.

\subsection{La categoria Nutrition}

Il lavoro si limita alla categoria \emph{Nutrition}. La categoria comprende 200 domande divise in tre prototipi, riassunti nella
Tabella~\ref{tab:prototipi}. Tutte le domande hanno cinque opzioni, per cui il
baseline casuale è del 20\%. I video coinvolti sono 50, appartenenti a 8
partecipanti diversi.

\begin{table}[htbp]
\centering
\small
\begin{tabular}{@{}lccl@{}}
\toprule
Prototipo & $N$ & Input & Domanda \\
\midrule
\texttt{video\_nutrition\_estimation} & 50 & 1 video & ingrediente con più \\
 & & & carbs/fat/protein/calorie \\[2pt]
\texttt{nutrition\_change} & 50 & 1 video + & variazione dei 4 valori \\
 & & intervallo & fra due istanti \\[2pt]
\texttt{image\_nutrition\_estimation} & 100 & 5 frame & quale dei 5 ingredienti \\
 & & (5 video) & mostrati ha più X \\
\bottomrule
\end{tabular}
\caption{I tre prototipi della categoria \emph{Nutrition}. Ogni domanda ha cinque opzioni.}
\label{tab:prototipi}
\end{table}

I tre prototipi non richiedono lo stesso tipo di ragionamento. In
\textbf{video\_nutrition\allowbreak\_estimation} la domanda riguarda l'intera ricetta e le
opzioni sono nomi di ingredienti; in \textbf{nutrition\allowbreak\_change} la domanda è
delimitata da due timestamp espliciti e le opzioni sono quattro valori
numerici; in \textbf{image\_nutrition\allowbreak\_estimation} l'input non è un video ma
cinque frame isolati, estratti da cinque video diversi a istanti indicati
nelle annotazioni. Per questa ultima differenza il prototipo resterà escluso dal secondo
esperimento.

Per esempio, una domanda di \textbf{video\_nutrition\allowbreak\_estimation} è
\emph{``What is the ingredient with highest carbs in this recipe?''}, con opzioni
\emph{butternut squash}, \emph{onions}, \emph{spring greens}, \emph{split red
lentils} e una quinta opzione che raggruppa più ingredienti; la risposta corretta
è la quarta.

\section{Preprocessing dei dati}

La preparazione dei dati è stata una parte importante e non trascurabile del lavoro.

\subsection{Download}

Gli identificativi dei video necessari sono stati estratti dai tre file di
annotazione della categoria, ottenendo 50 video distinti appartenenti a 8
partecipanti. Il download è avvenuto in due passaggi: prima tramite il
downloader ufficiale del dataset, poi con uno script dedicato che verifica quali
file risultino mancanti per ciascun prototipo e li recupera in parallelo dal
repository pubblico. Al termine tutti i 50 video sono stati ottenuti in forma
utilizzabile: quelli il cui primo download era incompleto sono stati scaricati
di nuovo e verificati con il checksum MD5 pubblicato insieme al dataset.

\subsection{Sottocampionamento}

I video originali durano in media circa 16 minuti, il che li
rende non facilmente utilizzabili
nell'ambiente di esecuzione. Ogni video è quindi stato ridotto a 128
frame, scelti in modo uniforme su tutta la sua durata e riscritti come
video a 1~fps alla risoluzione originale di $1408\times1408$ pixel. La copertura
temporale resta quella dell'intero video, mentre la densità di campionamento
scende: si tratta di un'approssimazione rispetto al protocollo originale.

\subsection{Estrazione dei frame per \texttt{image\_nutrition\_estimation}}

Il prototipo \textbf{image\_nutrition\allowbreak\_estimation} non può usare i video
sottocampionati: ogni domanda richiede cinque frame a timestamp esatti, mentre il
video a 128 frame contiene un frame ogni 14 secondi circa del video originale, e
il frame richiesto in generale non è fra questi. I 132 frame distinti richiesti
dalle 100 domande sono quindi stati estratti dai video integrali, all'istante
indicato nelle annotazioni. Per non dover scaricare per intero i 47 video
coinvolti (circa 58~GB), i frame sono stati letti direttamente dal repository
pubblico del dataset, scaricando solo la porzione di file attorno a ciascun
timestamp.

\subsection{Generazione dei pannelli}
\label{sec:panels-gen}

Per il secondo esperimento è stata reimplementata la procedura di creazione
dei Video Panels a partire dal codice ufficiale del Paper~\cite{videopanels}. Dato un video, si estraggono $C \cdot \alpha \cdot \beta$ frame
uniformi e li si dispone in $C$ griglie $\alpha \times \beta$ di frame
consecutivi fra quelli estratti. Sono stati usati i parametri indicati dagli autori del Paper come
configurazione migliore: $C=32$, $\alpha=\beta=2$, bordo nullo fra i riquadri e
stride temporale $\gamma=1$. Ogni video produce quindi 32 pannelli che coprono complessivamente 128 frame, gli stessi 128 del video
sottocampionato.  Viene quindi esteso il contesto temporale mentre diminuisce il contesto spaziale. La Figura~\ref{fig:panel} ne mostra un esempio.

Il metodo si applica solo quando il video è abbastanza lungo da rendere
sensato il processo; in caso contrario si ricade sul campionamento
uniforme ordinario.

Prima di ogni esecuzione viene inoltre eseguito un controllo che apre davvero
ciascun file richiesto e separa le domande eseguibili da quelle da escludere:
una domanda il cui input non è leggibile produrrebbe un'eccezione che verrebbe
contata come risposta sbagliata, abbassando l'accuracy. Con il dataset completo
tutte le 200 domande del primo esperimento risultano eseguibili.

\begin{figure}[htbp]
    \centering
    \includegraphics[width=0.5\linewidth]{figure.jpg}
    \caption{Un pannello $2\times2$: quattro frame consecutivi del video
occupano lo spazio di un singolo frame in ingresso al modello.}
    \label{fig:panel}
\end{figure}

\section{Valutazione e Metriche}
\label{sec:protocollo}

Affinché il confronto risulti sensato, tutto tranne il modello è stato gestito in maniera identica. Prompt, interpretazione
delle risposte e calcolo delle metriche sono quindi raccolti in un unico modulo
condiviso dai due esperimenti sui due modelli: l'unica parte specifica del modello è una funzione
di inferenza che riceve l'input visivo e il prompt e restituisce il testo
generato. Se ciascun modello avesse il proprio parser, la differenza di
accuracy misurata comprenderebbe anche la differenza fra i parser.

\subsection{Prompt}

Il prompt è deliberatamente minimale e riproduce il formato del paper: la
domanda, le cinque opzioni etichettate da A a E, e l'istruzione finale
\emph{``Answer with the option's letter from the given choices directly.''}.

Non viene usato alcun system prompt, perché il paper di riferimento non ne
prevede uno.

Viene impostato in fase di richiesta un
limite di 32 token in uscita. La risposta attesa è una singola lettera, quindi il
limite non tronca risposte legittime.

\subsection{Interpretazione delle risposte}

Un modello non sempre risponde nel formato richiesto: può anteporre
punteggiatura, premettere una formula come \emph{``the answer is''}, ricopiare il
testo dell'opzione invece della lettera, o rispondere con una cifra. Il parsing è
quindi organizzato in cinque livelli, applicati in ordine dal più stretto al più
permissivo:

\begin{enumerate}\itemsep2pt
  \item lettera all'inizio, eventualmente fra parentesi o seguita da
        punteggiatura;
  \item formula esplicita del tipo \emph{answer/option/choice} seguita da una
        lettera;
  \item testo di una delle opzioni ricopiato, purché corrisponda a una sola
        opzione;
  \item una sola lettera valida isolata nel resto del testo;
  \item una cifra da 1 a 5, interpretata come indice dell'opzione.
\end{enumerate}

Se nessun livello si applica, la risposta viene marcata come non interpretabile e
conteggiata come errata. Il livello che ha prodotto l'esito viene salvato insieme
alla risposta grezza del modello, così da poter sempre ispezionare \emph{come} si
è arrivati a un certo indice. Il primo livello è costruito in modo da non
confondere l'iniziale di una parola con un'etichetta, evitando per esempio di
leggere come opzione~C la \emph{C} di \emph{``Calories changed by\ldots''}.

\subsection{Metriche}

Il paper di riferimento definisce l'accuracy di una categoria come la
\textbf{media delle accuracy dei singoli prototipi}, non come l'accuracy
aggregata su tutte le domande. Con prototipi di dimensione diversa (50, 50 e 100
domande) i due valori non coincidono, per cui entrambi vengono riportati: il
primo è quello confrontabile con i valori pubblicati, il secondo è indicato come
accuracy aggregata. Il riferimento casuale è 20\%.

Oltre all'accuracy vengono riportate due diagnostiche. La prima è la
distribuzione dei metodi di parsing, che dice quanta accuracy si stia
eventualmente perdendo per ragioni di formato anziché di contenuto. La seconda è
il \textbf{bias posizionale}: la distribuzione delle lettere scelte dal modello
messa a confronto con quella delle risposte corrette. Se un modello fosse
soltanto impreciso, le due distribuzioni si somiglierebbero; se invece sceglie sempre le stesse lettere mentre la risposta corretta è dispersa, non sta
rispondendo alla domanda ma seguendo una preferenza a priori sul formato. È la
distinzione fra \emph{il modello sbaglia} e \emph{il modello non sta guardando il
video}, ed è la spiegazione tipica di un'accuracy inferiore al caso.

\subsection{Esecuzione}

Ogni risposta viene salvata insieme
alla risposta grezza del modello, all'indice predetto, a quello corretto e al metodo di
parsing: questo rende ogni esecuzione riprendibile dopo un'interruzione
dell'ambiente e, soprattutto, rende i numeri di accuracy verificabili a
posteriori riga per riga. Le risposte terminate con un errore di esecuzione
(per esempio memoria GPU esaurita) non vengono considerate completate e sono
ripetute alla ripresa successiva.

\section{Esperimento 1: Valutazione di VideoLLaMA 2 e LLaVA-Video}
\label{sec:baseline}

\subsection{Setup}

Entrambi i modelli sono stati eseguiti su una singola GPU NVIDIA A100 da 80~GB.

\textbf{VideoLLaMA~2 7B}~\cite{videollama2} è stato caricato dal checkpoint
ufficiale in precisione \texttt{float16}. Il repository forza l'uso di
FlashAttention sulla vision tower, che però non lo supporta: è stato necessario
sostituirla con l'implementazione \emph{eager} per rendere il modello
eseguibile. L'input visivo passa sempre dal preprocessore del repository, che
applica ridimensionamento a $336\times336$, normalizzazione CLIP e padding, sia
che riceva il percorso di un video sia che riceva una lista di immagini. Un
dettaglio rilevante è che il preprocessore ricampiona l'input a
\texttt{num\_frames}, il cui valore di default è \textbf{8}: dei 128 frame
preparati a monte, il modello ne vede effettivamente otto.

\textbf{LLaVA-Video 7B}~\cite{llavavideo} è stato caricato tramite il loader
ufficiale in precisione \texttt{bfloat16} con attenzione \texttt{sdpa}. A
differenza del precedente, richiede frame già estratti: il campionamento
uniforme viene fatto esplicitamente a \textbf{32}
frame. La scelta verrà motivata nella sezione riguardante l'esperimento dei Panels.

I due modelli differiscono quindi per quanti frame ricevono. È una
differenza che va tenuta presente nel confronto e su cui si torna fra i limiti.

\subsection{Risultati}

La Tabella~\ref{tab:baseline} confronta i valori ottenuti con quelli pubblicati.
Entrambe le baseline sono riprodotte: VideoLLaMA~2 risulta 4.0 punti sopra il
valore di riferimento, mentre LLaVA-Video coincide con il valore pubblicato. L'ordinamento fra i due
modelli è rispettato, e in entrambi i casi il valore si colloca nettamente sopra
il 20\% del caso e altrettanto nettamente sotto l'85\% degli annotatori umani.

\begin{table}[htbp]
\centering
\small
\begin{tabular}{@{}lccc@{}}
\toprule
Modello & Paper & Riprodotto & $\Delta$ \\
\midrule
VideoLLaMA 2 7B & 32.7 & 36.7 & $+4.0$ \\
LLaVA-Video 7B  & 38.7 & 38.7 & $+0.0$ \\
\midrule
Umano           & 85.0 & --- & \\
Casuale         & 20.0 & --- & \\
\bottomrule
\end{tabular}
\caption{accuracy sulla categoria \emph{Nutrition} (\%), come media delle
accuracy per prototipo.}
\label{tab:baseline}
\end{table}

Il sottocampionamento dei video rende il nostro protocollo diverso da quello
originale, e una differenza di qualche punto è attesa: con 50 domande per
prototipo una sola risposta vale 2 punti. Il risultato da trattenere è che i
valori sono compatibili con quelli pubblicati, non che siano superiori.

\subsection{Dettaglio per prototipo}

L'accuracy aggregata nasconde un comportamento molto disomogeneo, riportato
nella Tabella~\ref{tab:per-prototipo}.

\begin{table}[htbp]
\centering
\small
\begin{tabular}{@{}lcc@{}}
\toprule
Prototipo ($N$) & VideoLLaMA 2 & LLaVA-Video \\
\midrule
\texttt{video\_nutrition\_estimation} (50) & 66.0 & 68.0 \\
\texttt{nutrition\_change} (50) & 16.0 & 22.0 \\
\texttt{image\_nutrition\_estimation} (100) & 28.0 & 26.0 \\
\midrule
Media per prototipo             & 36.7 & 38.7 \\
Aggregata sulle domande         & 34.5 & 35.5 \\
\bottomrule
\end{tabular}
\caption{accuracy per prototipo (\%). I due modelli si comportano in modo
molto diverso a seconda del prototipo.}
\label{tab:per-prototipo}
\end{table}

Tre osservazioni:

\textbf{Il prototipo \texttt{nutrition\_change} è al livello casuale o sotto.}
VideoLLaMA~2 ottiene 16.0\%, cioè meno del 20\% casuale. Il bias posizionale
infatti chiarisce: il modello colloca il 58\% delle proprie
risposte sull'opzione B, mentre le risposte corrette sono distribuite quasi
uniformemente sulle cinque lettere. LLaVA-Video mostra lo stesso fenomeno in
forma attenuata, con il 42\% delle risposte su B e quasi nessuna su E. Su questo
prototipo i modelli non stanno rispondendo alla domanda: stanno seguendo una
preferenza a priori sul formato. È coerente con la natura del compito, che
richiede di leggere valori numerici su una bilancia digitale in un intervallo di
pochi secondi e di sceglierne la variazione fra cinque quaterne di numeri simili.
A questo si aggiunge che il modello riceve frame campionati uniformemente su
tutto il video senza l'indicazione dell'istante a cui ciascuno corrisponde:
i timestamp citati nella domanda non sono quindi collegabili ai frame, e
un intervallo di pochi secondi in un video di un quarto d'ora quasi sempre non
cade su nessuno dei frame ricevuti.

\textbf{Il prototipo \texttt{video\_nutrition\allowbreak\_estimation} è molto sopra la media
di categoria}, con il 66\% per VideoLLaMA~2 e il 68\% per LLaVA-Video. Una spiegazione
plausibile è che molte di queste domande siano risolvibili con la sola conoscenza
linguistica, senza guardare il video: sapere che le lenticchie contengono più
carboidrati della cipolla non richiede evidenza visiva.

\textbf{Il prototipo \texttt{image\_nutrition\allowbreak\_estimation} resta vicino al
caso}, con il 28\% e il 26\%. Ogni domanda chiede quale dei cinque ingredienti
inquadrati abbia più calorie, grassi, carboidrati o proteine, e anche qui le
opzioni sono nomi di ingredienti. Il divario con il prototipo precedente
suggerisce che la sola conoscenza linguistica non basti: un'ipotesi plausibile,
non verificata in questo lavoro, è che i valori nutrizionali delle annotazioni
si riferiscano alla quantità effettivamente usata nella ricetta, visibile solo
nelle immagini, e non all'ingrediente in astratto.

\section{Esperimento 2: Valutazione dei VLM con la tecnica dei Video Panels}
\label{sec:panels}

\subsection{Il metodo e la condizione del confronto}

I Video Panels~\cite{videopanels} affrontano un vincolo strutturale dei VLM: la
finestra di contesto limita quanti frame possono essere dati in pasto al
modello, e su un video lungo questo significa un campionamento molto rado. Si può quindi ridurre risoluzione spaziale dove non serve: ogni elemento
della sequenza in ingresso, invece di contenere un frame, contiene una
griglia $\alpha \times \beta$ di frame consecutivi. A parità di elementi in
ingresso la copertura temporale si moltiplica per $\alpha\beta$, riducendo di altrettanto il contesto spaziale di ciascun frame.

Il confronto è quindi sensato solo \emph{a parità di finestra di contesto}, ed è
in questa condizione che è stato condotto. Per VideoLLaMA~2, che ricampiona
internamente a 8 elementi, si confrontano 8 frame contro 8 pannelli, cioè 8
contro 32 frame di copertura. Per LLaVA-Video, che lavora con 32 frame, si
confrontano 32 frame contro 32 pannelli, cioè 32 contro 128. In entrambi i
casi il modello riceve lo stesso numero di elementi visivi nelle due condizioni e
quadruplica la copertura temporale passando ai pannelli; cambia il regime
assoluto in cui la cosa avviene.

L'esperimento è limitato ai due prototipi basati su video, per un totale di 99
domande valutate. \textbf{image\_nutrition\allowbreak\_estimation} è escluso in quanto il suo input è costituito da cinque frame isolati provenienti
da cinque video diversi, quindi non esiste alcuna sequenza temporale da
impacchettare e il metodo non è applicabile.

\subsection{Risultati}

\begin{table}[htbp]
\centering
\small
\begin{tabular}{@{}lcccc@{}}
\toprule
& \multicolumn{2}{c}{VideoLLaMA 2} & \multicolumn{2}{c}{LLaVA-Video} \\
\cmidrule(lr){2-3}\cmidrule(lr){4-5}
Prototipo & base & panels & base & panels \\
\midrule
\texttt{video\_nutrition\_estimation} & 66.0 & 62.0 & 68.0 & 72.0 \\
\texttt{nutrition\_change} & 16.3 & 20.4 & 20.4 & 14.3 \\
\midrule
Media per prototipo        & 41.2 & 41.2 & 44.2 & 43.2 \\
$\Delta$                   & \multicolumn{2}{c}{$0.0$} & \multicolumn{2}{c}{$-1.0$} \\
\bottomrule
\end{tabular}
\caption{Video liscio contro pannelli $2\times2$, a parità di finestra di
contesto e sulle stesse 99 domande. I valori del baseline differiscono da quelli
della Tabella~\ref{tab:per-prototipo} perché calcolati sul sottoinsieme
ristretto ai due prototipi video.}
\label{tab:panels}
\end{table}

Il risultato aggregato, riportato nella Tabella~\ref{tab:panels}, è che
\textbf{i Video Panels non producono alcun guadagno} sulla categoria
\emph{Nutrition}: $0.0$ punti per VideoLLaMA~2 e $-1.0$ per LLaVA-Video.
Entrambe le variazioni sono troppo piccole per essere distinte dal rumore di
campionamento e vanno lette come assenza di effetto, non come peggioramento.

È interessante notare che i due modelli
si spostano in \textbf{direzioni opposte sugli stessi prototipi}, con variazioni
di ampiezza simile e segno contrario: VideoLLaMA~2 guadagna 4.1 punti su
\texttt{nutrition\_change} e ne perde 4.0 su
\texttt{video\_nutrition\_estimation}, mentre LLaVA-Video fa il contrario,
guadagnando 4.0 punti su \texttt{video\_nutrition\_estimation} e perdendone 6.1 su
\texttt{nutrition\_change}.

Va detto con chiarezza che su prototipi da circa 50 domande una variazione di 4
punti corrisponde a due domande: nessuno di questi scarti è statisticamente
significativo preso singolarmente.

\subsection{Osservazioni sul Visual Prompting}

Applicare la tecnica dei Video Panels (\textbf{Visual Prompting}) non ha portato i vantaggi desiderati. I Video Panels scambiano risoluzione spaziale per risoluzione
temporale, e questo è vantaggioso quando l'informazione rilevante è
\emph{dove e quando} accade qualcosa di visivamente evidente. Il prototipo
\texttt{nutrition\_change} richiede però di leggere \emph{cifre} sul display di
una bilancia digitale: dividendo per quattro l'area di ciascun frame si
degrada esattamente il tipo di dettaglio da cui dipende la risposta. Questo è
coerente con il calo di LLaVA-Video su questo prototipo; VideoLLaMA~2, che nel
baseline vede solo 8 frame, mostra invece un lieve miglioramento, compatibile con
un maggior beneficio della copertura temporale quando i frame sono molto pochi.
In entrambi i casi, però, le variazioni restano entro il rumore di campionamento
e i due modelli rimangono vicini al livello casuale.

\section{Discussione e conclusioni}

\subsection{Limiti}

I risultati vanno letti tenendo presenti alcuni limiti, che riguardano sia i dati
sia il disegno sperimentale.

\textbf{Dati.} I video sono stati sottocampionati a 128 fotogrammi anziché usati
per intero, per i due prototipi basati su video; solo
\texttt{image\_nutrition\_estimation} usa i frame originali agli istanti
indicati dalle annotazioni. L'insieme di valutazione coincide invece con quello
originale: tutte le 200 domande sono state valutate.

\textbf{Frame ricevuti.} I due modelli vedono un numero diverso di frame (8 per
VideoLLaMA~2, 32 per LLaVA-Video), per cui il confronto diretto fra i due
modelli mescola l'effetto dell'architettura con quello della quantità di input.

\textbf{Riferimento temporale.} Le domande di \texttt{nutrition\_change} citano
due istanti precisi del video, ma i modelli ricevono frame campionati
uniformemente senza l'indicazione del tempo corrispondente; in questo protocollo
il prototipo misura quindi anche la capacità di collegare i timestamp della
domanda a frame che quasi sempre non li contengono.

\textbf{Ambito.} Il lavoro copre una sola delle sette categorie del benchmark.
Le conclusioni valgono per \emph{Nutrition} e non sono estendibili alle altre
senza ulteriori esperimenti.

\subsection{Sviluppi possibili}

Si individuano tre direzioni. La prima è
eseguire un \textbf{baseline testuale}, cioè sottoporre le stesse domande a un
modello senza accesso al video: è la misura che manca per stabilire quanta parte
dell'accuracy su \texttt{video\_nutrition\_estimation} dipenda da conoscenza
linguistica anziché da comprensione visiva, ipotesi che qui resta formulata ma
non verificata. La seconda è variare la geometria dei pannelli, in particolare
verso griglie più grandi su un modello con finestra di contesto ampia, per capire
se il compromesso fra risoluzione spaziale e temporale abbia un punto ottimale
diverso su questo dominio. La terza è applicare i panels non all'intero video
ma solo alle parti del video più rilevanti per rispondere alla domanda, per
esempio l'intervallo indicato nelle domande di \texttt{nutrition\_change}.

\subsection{Conclusioni}

Il lavoro ha riprodotto le baseline della categoria \emph{Nutrition} del
benchmark VQA di HD-EPIC~\cite{hdepic} per due VLM open da 7 miliardi di parametri, ottenendo
36.7\% per VideoLLaMA~2~\cite{videollama2} e 38.7\% per LLaVA-Video~\cite{llavavideo}, contro i 32.7\% e 38.7\%
riportati dagli autori. Su questa base ha poi valutato
i Video Panels~\cite{videopanels} come metodo di visual prompting, trovando che a parità di finestra
di contesto non producono alcun guadagno sulla categoria, e che il loro effetto
si distribuisce in modo opposto sui due prototipi a seconda del modello.

\begin{thebibliography}{9}

\bibitem{hdepic}
T.~Perrett et al.
\newblock \emph{HD-EPIC: A Highly-Detailed Egocentric Video Dataset}.
\newblock In \emph{Proceedings of the IEEE/CVF Conference on Computer Vision and
Pattern Recognition (CVPR)}, 2025.

\bibitem{videollama2}
Z.~Cheng et al.
\newblock \emph{VideoLLaMA 2: Advancing Spatial-Temporal Modeling and Audio
Understanding in Video-LLMs}.
\newblock arXiv:2406.07476, 2024.

\bibitem{llavavideo}
Y.~Zhang et al.
\newblock \emph{Video Instruction Tuning with Synthetic Data}.
\newblock arXiv:2410.02713, 2024.

\bibitem{videopanels}
L.~Doorenbos, F.~Spurio, J.~Gall.
\newblock \emph{Video Panels for Long Video Understanding}.
\newblock arXiv:2509.23724, 2025.

\end{thebibliography}

\end{document}
